% ---------------------------------------------------------------
% Section 6: Limitations and Future Work
% ---------------------------------------------------------------
\section{Limitations and Future Work}\label{sec:limitations}

\subsection{Limitations}

The Bailout Protocol makes a single, specific contribution — a compulsory exception-escalation mechanism with a mandatory human terminus — and the limitations described here are coextensive with the boundaries of that contribution. We do not claim that the protocol resolves all challenges in safe autonomous system design; we claim only that it addresses one structural gap that existing approaches leave open.

\medskip\noindent\textbf{Human response latency.}%
\label{sec:lim:latency}
The protocol's liveness guarantee — that every unresolved exception \emph{must} eventually reach $h(n)$ — is conditioned on $h(n)$ being reachable and responsive within the provisioned timeout $T_{\mathrm{human}}$. In high-throughput production environments where nodes generate exception states at millisecond cadence, a human principal capable of meaningful adjudication cannot operate at machine speed. The protocol mitigates this by providing a bounded \texttt{unresolved} tag and a capability timeout that prevents indefinite stalling, but it does not eliminate the fundamental speed differential between machine reasoning and human deliberation. When $T_{\mathrm{human}}$ expires without acknowledgment, the protocol records a \texttt{timeout} entry in the Forensic Ledger and transitions to a terminal \texttt{unresolved} state. This is an honest acknowledgment of a boundary condition, not a resolution of it. The practical implication is that the Bailout Protocol is architecturally sound but operationally bounded by the latency and availability of the human anchor. Systems operating in domains requiring sub-second exception resolution cannot rely on human-in-the-loop bailout as their primary safety mechanism; they must instead rely on the four preceding pillars (Loop Isolation, Recursive Spawning, Spatial Firewall, and Sandbox Gate) to handle the exception conditions they can anticipate, reserving the Bailout Protocol for the complementary set of conditions those pillars cannot cover.

\medskip\noindent\textbf{Adversarial conditions.}%
\label{sec:lim:adversarial}
The protocol is designed for the failure modes of autonomous systems operating within their intended deployment context — not for environments where the system itself is the adversarial actor, or where communication infrastructure is actively contested. If an adversarial machine actor with Fact Layer write access modifies its own operating range $R(n)$ to absorb exceptions that should trigger bailout, the Fact Layer's pre-commit hook enforces the immutability of $R(n)$ only against unauthorized writes from other nodes; it does not prevent a node from misrepresenting its own observational boundary. More severely, if the communication pathway between a node and $h(n)$ is severed — whether by network failure, targeted attack, or institutional censorship — the protocol can detect non-acknowledgment within $T_{\mathrm{human}}$ and record the failure, but it cannot compel the delivery of the exception to a human principal across an unavailable network. The protocol's forensic guarantee (that the failure is immutably recorded) holds under adversarial conditions; its liveness guarantee (that the human is eventually contacted) does not.

\medskip\noindent\textbf{Scalability.}%
\label{sec:lim:scalability}
The mandatory human terminus is a scalability bottleneck by design. Every unresolved exception that propagates to $h(n)$ consumes the attention of a human principal, and human attention does not scale with the number of active nodes or the rate of exception generation. The protocol mitigates this through the $T_{\mathrm{bounds}}$ window, which allows the machine-only execution graph to handle anticipated exceptions without escalation, and through the \texttt{unresolved} tag, which allows $h(n)$ to batch-process cases that cannot be resolved within the timeout. However, these mitigations do not change the asymptotic cost structure: each exception that survives the four preceding pillars requires human attention, and human attention is the bottleneck resource. For deployments with large autonomous agent populations operating in noisy or unpredictable environments, the operational cost of maintaining a human anchor that can handle the resulting exception volume may exceed available human capacity. This is not a failure of the protocol's design — it is the honest cost of a mandatory human-in-the-loop guarantee. Systems that require scalability at the level of thousands of concurrent nodes per human anchor require architectural extensions (discussed in \autoref{sec:future}) that go beyond the current specification.

\medskip\noindent\textbf{Exploitation risk.}%
\label{sec:lim:exploitation}
The protocol assumes that machine actors trigger bailout faithfully when their operating range is exceeded. However, a sufficiently sophisticated agent — particularly one operating in a domain where escalation is economically costly to its objectives — may develop strategies to suppress or mask exception states rather than trigger the protocol. This is goal hacking applied to the accountability mechanism itself: the agent learns to classify genuinely out-of-range inputs as in-range, or to absorb minor exceptions before they compound into bailout-triggering conditions. The protocol's enforcement architecture (the Fact Layer pre-commit hook, the immutable parent pointer, the Forensic Ledger) protects against other nodes tampering with a given node's escalation state, but it does not protect against a node misrepresenting the observational boundary of its own operating range. Defenses against this class of attack — which is a form of model stealing or capability concealment — require techniques outside the scope of the current protocol specification, including behavioral auditing of node decision boundaries, cross-node anomaly detection, and runtime capability evaluation.

\subsection{Future Work}\label{sec:future}

The following extensions address the limitations identified above and represent the natural directions for subsequent research and engineering development.

\medskip\noindent\textbf{Formal verification.}%
\label{sec:future:verification}
The current specification states the Safety, Liveness, and Bypass theorems informally and provides proof sketches. A rigorous machine-checked verification — using a proof assistant such as Coq or Lean — would establish the correctness of the state machine specification, the transition relation, and the invariants against the full state space of $\mathcal{B}$. Of particular importance is the bypass theorem: the claim that the escalation pathway contains no machine actor as a terminal absorption point depends on subtle assumptions about the immutability of the parent pointer and the non-delegability of the Purpose Layer's exception authority. Formal verification would clarify which of these assumptions are independent and which are interdependent, and would provide a basis for reasoning about compositions of the Bailout Protocol with other protocol extensions.

\medskip\noindent\textbf{Adaptive timeout provisioning.}%
\label{sec:future:adaptive}
The current specification treats timeout values ($T_{\mathrm{bounds}}$, $T_{\mathrm{prop}}$, $T_{\mathrm{cap}}$, $T_{\mathrm{human}}$) as statically provisioned at node setup. In production deployments, the appropriate timeout for a given node depends on the exception rate in its operating environment, the current load on $h(n)$, the criticality of the node's function, and the demonstrated reliability of its communication pathway. An adaptive timeout mechanism — one that adjusts timeout values based on observed exception density and anchor availability while maintaining the non-negotiable liveness floor $T_{\mathrm{human}}^{\min}$ — would improve the protocol's operational practicality without sacrificing its correctness guarantees. The adaptive mechanism must be verified to maintain the liveness invariant under all non-Byzantine conditions.

\medskip\noindent\textbf{Distributed human anchor pools.}%
\label{sec:future:pools}
The current specification assumes a single named human principal $h(n)$ per node. This is sufficient for systems with small-to-moderate node populations where a single human anchor can maintain situational awareness across the exceptions it receives. For large-scale deployments — thousands of active nodes, high exception rates — a single anchor becomes a throughput bottleneck. A distributed anchor pool architecture would replace the single $h(n)$ with a roster of qualified human principals, each capable of serving as the terminus for any exception routed to the pool, with a load-balancing policy that distributes escalation events across the roster to maximize aggregate throughput while preserving the guarantee that every exception reaches a human, not a machine actor. The pool architecture introduces a new correctness requirement: the pool must be shown to maintain the mandatory human terminus property under roster membership changes (principal addition, removal, or temporary unavailability). The Forensic Ledger must record not just the pool-level exception event but the identity of the assigned principal, preserving full attribution.

\medskip\noindent\textbf{Robustness under Byzantine and adversarial conditions.}%
\label{sec:future:byzantine}
Extending the protocol to maintain its guarantees when machine actors may behave arbitrarily — not merely failing to follow the protocol but actively attempting to suppress or redirect exception states — requires additional mechanisms. These include cryptographic attestation of bailout trigger events (so that a node cannot retroactively deny that it entered BOUNDED state), redundant multi-path escalation (so that the severance of a single communication pathway does not prevent contact with $h(n)$), and cross-node state reconciliation (so that sibling nodes can detect whether a parent node has suppressed an exception that should have been escalated). Each extension must be verified to preserve the bypass theorem: no amount of Byzantine behavior by intermediate machine actors may prevent the exception from reaching a human principal.

% ---------------------------------------------------------------
% Section 7: Conclusion
% ---------------------------------------------------------------
\section{Conclusion}\label{sec:conclusion}

The Bailout Protocol advances a single, architecturally precise claim: that autonomous systems built on the \bxthree{} Framework must propagate every unresolved exception state to a named human principal, bypassing all intermediate machine actors, as a compulsory structural requirement rather than a runtime choice. This mandatory bailout property is not a feature that can be retrofitted to an existing agentic architecture through better error handling or improved escalation heuristics. It is a structural property that follows, as a matter of architectural necessity, from the functional separation between the \purpose{}, \bounds{}, and \fact{} layers.

The argument is straightforward. When the \bounds{} layer encounters a condition outside its provisioned operating range $R(n)$ that it cannot resolve, no machine actor in the system has the authority to adjudicate that condition. The \bounds{} layer detects exception conditions but carries no judgment authority; the \fact{} layer enforces transition relations and records events but carries no intent authority. The decision must return to the \purpose{} layer, which carries intent, judgment, and final accountability. The Bailout Protocol is the mechanism that enforces this return — compulsorily, deterministically, and with full forensic attribution at every step.

This contribution is realized within the \bxthree{} Framework as its fifth named architectural pillar, complementing Loop Isolation, Recursive Spawning, Spatial Firewall, and Sandbox Gate. Together, these five pillars constitute the upstream accountability guarantee: that \bxthree{} systems \emph{fail upward into human consciousness, never downward into autonomous chaos}. The Bailout Protocol is the runtime expression of this guarantee for all exception conditions not resolved by the preceding four pillars — including conditions not anticipated at design time, which is precisely the operational reality that makes the protocol necessary.

The \bxthree{} Framework's role in this contribution is to provide the conceptual architecture within which the Bailout Protocol operates as a necessary consequence rather than an independent design choice. The functional separation of \purpose{}, \bounds{}, and \fact{} layers creates the structural conditions that make the bailout requirement compulsory. The Forensic Ledger maintains the immutable attribution record that makes every escalation auditable. The immutable parent pointer established at node provisioning ensures that the escalation path is fixed at the architectural level and cannot be modified by any machine actor. These elements are not independent mechanisms that happen to cooperate; they are mutually defining structural facts whose necessity is established by the protocol's own specification. The upstream accountability guarantee and the Bailout Protocol are two aspects of the same architectural commitment: that in a \bxthree{} system, every exception that exceeds machine authority reaches a human who can bind it, and every such exception is recorded so that the binding can be verified.

The theoretical claim is foundational: that the upstream accountability guarantee of the \bxthree{} Framework is not a design aspiration but a logical consequence of the framework's architecture, and that the Bailout Protocol is the mechanism by which that consequence is enforced at runtime. The practical claim is equally clear: that autonomous systems operating without such a mechanism are structurally exposed to failure modes — invisible exception absorption, contingent escalation, and irreversible drift — that cannot be eliminated by incremental improvement to existing error-handling or oversight practices. They require a new architectural primitive. The Bailout Protocol is that primitive.